\section*{5. Appendix: Integrality Theorems}

\begin{center}
By P. Deligne
\end{center}

\subsection*{5.0.}
Let $p$ be a prime number, $q$ a power of $p$, $\mathbb F_q$ a field with
$q$ elements, $E$ the completion of a number field at a place dividing a
prime number $\ell\ne p$, and $T$ a set of prime numbers (for example,
$T=\varnothing$).

An element $\alpha$ of an algebraic closure of $E$ is called
\emph{$T$-integral} if it is algebraic over $\mathbb Q$ and integral over
$\mathbb Z[(1/r)_{r\in T}]$.  An endomorphism $F$ of a finite-dimensional
$E$-vector space is called $T$-integral if its eigenvalues are $T$-integral.

Let $k$ be a finite field with $q'$ elements, let $\bar k$ be an algebraic
closure of $k$, and let
\[
\mathfrak G=\operatorname{Gal}(\bar k/k)
\]
be the Galois group of $k$.  We denote by $f\in\mathfrak G$ the Frobenius
substitution
\[
f(x)=x^{q'},
\]
and by $F$ the ``geometric'' Frobenius, the inverse of $f$ in $\mathfrak G$.
A representation $\rho$ of $\mathfrak G$ on a finite-dimensional
$E$-vector space is called $T$-integral if $\rho(F)$ is $T$-integral.

If $X$ is a scheme of finite type over $\mathbb F_q$ and $x$ is a closed
point of $X$, we denote by $F_x$ the geometric Frobenius element of the
Galois group of $k(x)$.  If $\bar x$ is a geometric point of $X$ lying over
$x$ and $\underline G$ is an $E$-sheaf on $X$, then $F_x$ acts, by transport
of structure, on the fiber $G_{\bar x}$.

\subsubsection*{Definition 5.1.}
A constructible $E$-sheaf $\underline G$ on a scheme $X$ of finite type over
$\mathbb F_q$ is called $T$-integral if, for every geometric point $\bar x$
of $X$ lying over a closed point $x$ of $X$, $F_x$ is a $T$-integral
automorphism of $G_{\bar x}$.

If $a:X\longrightarrow\operatorname{Spec}(\mathbb F_q)$ is a scheme of
finite type over $\mathbb F_q$ and $\underline G$ is a constructible
$E$-sheaf on $X$, we write $\underline H^*(X,\underline G)$ for the
$E$-sheaves $R^*a_*\underline G$ on
$\operatorname{Spec}(\mathbb F_q)_{\acute e t}$.  If
$\overline{\mathbb F}_q$ is an algebraic closure of $\mathbb F_q$ and
$\overline X=X\times_{\mathbb F_q}\overline{\mathbb F}_q$, these sheaves
identify with the cohomology groups $H^*(\overline X,\underline G)$ regarded
as Galois modules, that is, endowed with an action of the geometric
Frobenius $F$.  We use the same convention for cohomology with proper
support.

\subsubsection*{Lemma 5.2.1.}
Let $\underline G$ be a constructible $E$-sheaf on a scheme $X$ of finite
type over $\mathbb F_q$, of dimension at most $n$.

\begin{enumerate}[label=(\roman*)]
\item There is a closed subset $Y$ of $X$, of dimension $0$, such that the
  evident map
  \[
  \underline H^0(X,\underline G)\longrightarrow
  \underline H^0(Y,\underline G)
  \]
  is injective.
\item If $X$ is separated and $U$ is an open subset of $X$ whose complement
  $Z$ has dimension less than $n$, there is a closed subset $Y$ of $U$, of
  dimension $0$, and a surjective map
  \[
  \underline H^0(Y,\underline G)(-n)\longrightarrow
  \underline H_c^{2n}(X,\underline G).
  \]
\end{enumerate}

\medskip
\noindent\emph{Proof.}
Assertion (i) is evident.  We prove (ii).  After replacing $X$ by
$X_{\mathrm{red}}$ and shrinking $U$, we may suppose that $U$ is smooth and
that $\underline G|_U$ is twisted constant.  Since $\dim Z<n$, we have
\[
0=\underline H_c^{2n-1}(Z,\underline G)
\longrightarrow \underline H_c^{2n}(U,\underline G)
\xrightarrow{\sim}\underline H_c^{2n}(X,\underline G)
\longrightarrow \underline H_c^{2n}(Z,\underline G)=0,
\]
and Poincar\'e duality (SGA~4, Expos\'e~XVIII, and [2]) gives
\[
\underline H_c^{2n}(U,\underline G)\simeq
\underline H^0(U,\underline G^\vee(n))^\vee.
\]
The conclusion follows from (i).
